% CLAUDE REPORT::WORSTCASE-BRIDGE-0001::2026-11-12T13:00-08:00
\documentclass[11pt]{article}
\usepackage{amsmath,amssymb,amsthm,booktabs,longtable}
\usepackage[margin=1in]{geometry}
\newtheorem{theorem}{Theorem}
\newtheorem{lemma}[theorem]{Lemma}
\newtheorem{proposition}[theorem]{Proposition}
\newtheorem{corollary}[theorem]{Corollary}
\newcommand{\SAT}{\mathsf{SAT}}
\newcommand{\Pc}{\mathsf{P}}
\newcommand{\NP}{\mathsf{NP}}
\newcommand{\coNP}{\mathsf{coNP}}
\newcommand{\Wb}{\overline W}
\title{Worst-case derivation cost $\Wb_R(N)$ and the bridge $\SAT\in\Pc\Rightarrow\Wb_R\le\mathrm{poly}$:\\
a locality dichotomy for multisign hypergraph reduction systems}
\author{Claude (scheduled routine), run WORSTCASE-BRIDGE-0001}
\date{2026-11-12}
\begin{document}
\maketitle

\section*{Status and target}
Target (shared directive):
\[
\Wb_R(N)=\max_{\substack{F\in\mathrm{UNSAT}\\ |\langle F\rangle|\le N}}\ \min_{\pi:F\vdash_R\bot}\operatorname{cost}_R(\pi),
\qquad
\mathrm{Bridge}(R):\ \SAT\in\Pc\Longrightarrow \Wb_R(N)\le\mathrm{poly}(N).
\]
Results of this run: Lemma~\ref{lem:local} and Theorem~\ref{thm:dichotomy}
are \textbf{PROVED/EXACT} (elementary; the ingredients are known theorems and
the combination is stated here for the user's multisign hypergraph model).
Proposition~\ref{prop:refl} is a \textbf{KNOWN THEOREM} transfer.
$\mathrm{Bridge}(R)$ for every system with a known superpolynomial lower bound
is shown to be \emph{equivalent} to $\Pc\ne\NP$. $\mathrm{Bridge}(\mathrm{EF})$ and
superpolynomial lower bounds on $\Wb_{\mathrm{EF}}$ are \textbf{OPEN}.
\textbf{P versus NP remains OPEN.}

\section*{Exact quantifiers and model}
A \emph{multisign DAG hypergraph reduction system} $R$ is a Cook--Reckhow
proof system whose lines are signed hyperedges (clauses) over the variables of
$F$ and possibly auxiliary nodes, whose steps derive one hyperedge from at
most $k$ earlier ones by a polynomial-time checkable sound rule, and whose
derivations are DAGs (lines may be reused). $\operatorname{cost}_R(\pi)$ is the
number of lines (or total bit length; the two agree up to a polynomial for the
systems below). The minimum over $\pi$ grants \emph{free knowledge of the best
reduction sequence}, as the directive requires: search cost is never charged.
\emph{Local} means: no auxiliary nodes and every line has width at most $w$,
with $k,w$ constants. \emph{Extended} means auxiliary nodes may be introduced
by definitions $z\leftrightarrow\varphi(\text{earlier nodes})$.

\section*{Corpus reviewed}
Root \texttt{main.tex} (status sections), \texttt{AGENTS.md}; the newest
coordination handoffs of 2026-11-12 (LSPV-1/598--599, LSPV-2/580--581,
LSPV-3/594--595, LSPV-4/567--568, LSPV-5/509--510), which work on the
ARC/FLMV native replay and certificate-charge spine;
\texttt{research/agents/codex/2026-11-11\_1203\_SAT\_POLYNOMIAL\_ASSUMPTION\_AUDIT.tex}
(\S``A weaker sufficient bridge''), \texttt{\ldots\_1331\_TRANSFER\_RECOVERY\_AND\_RATE\_AMPLIFICATION.tex},
\texttt{\ldots\_2017\_ADAPTIVE\_COLLISION\_AND\_CLOSURE\_BARRIERS.tex},
\texttt{\ldots\_2035\_COMPATIBILITY\_CORE\_AND\_DEGREE\_TRANSFER.tex};
the earlier Claude world ledger (\texttt{research/agents/claude/world-ledger/}).
Not every historical LSPV report was reread; the search for
\texttt{W\_R}/worst-case derivation matched the Codex reports above and LSPV-5/458--467,
which concern restricted DAG abstractions, not the bridge.

\section*{Duplication / contradiction audit}
Codex 1203 already states (its eq.~(12) discussion) that once a lower bound on
$R$ is proved, the bridge is logically equivalent to the separation, and that
Cook--Reckhow gives only an \emph{existential} p-bounded system under
$\NP=\coNP$. This report \textbf{adopts} that statement as Proposition~\ref{prop:equiv}
and does not claim it as new. New here: the locality simulation
(Lemma~\ref{lem:local}) that places every local multisign reduction system in
the ``lower bound known, bridge $\equiv$ P$\ne$NP'' branch, the resulting
dichotomy, and the reflection reformulation of the extended branch. No
contradiction with the current root \texttt{main.tex} was found: nothing here
touches the ARC/FLMV strict box or any specialist ownership region.

\section*{Primary literature (checked this run)}
Cook--Reckhow (1979): a p-bounded proof system exists iff $\NP=\coNP$.
Haken (1985), Urquhart (1987), Chv\'atal--Szemer\'edi (1988), Ben-Sasson--Wigderson
(2001): exponential resolution lower bounds (PHP, Tseitin on expanders,
random 3-CNF) and width--size relation. Buss (1987): polynomial Frege proofs of
PHP. Cook (1975): PV proofs translate to polynomial-size EF proofs.
Krajíček--Pudlák (1989): for $Q$ a proof system and $R\ge_p\mathrm{EF}$ suitably
closed, $R\ge_p Q$ iff $R$ has polynomial-size proofs of the reflection
tautologies $\mathrm{RFN}_Q$. Recent: Krajíček (arXiv:1905.12935) on consistency
of lower bounds with PV; Arteche et al.\ (arXiv:2506.16956) on the proof analysis
problem for EF; ECCC 2026/109 on reflection and implicit proof systems (only an
excerpt seen; not used to support any claim); Lu--Santhanam--Tzameret (2025),
AC$^0[p]$-Frege lower bounds. None supplies $\mathrm{Bridge}(\mathrm{EF})$
or a superpolynomial EF/Frege lower bound.

\section{Free best sequence = nondeterministic proof length}
\begin{proposition}[KNOWN THEOREM, restated]\label{prop:cr}
$\exists R$ with $\Wb_R(N)\le\mathrm{poly}(N)$ $\iff$ $\NP=\coNP$.
\end{proposition}
This is Cook--Reckhow with cost $=$ proof length. Consequences for the program:
(a) granting the optimal sequence for free makes $\Wb_R$ a
\emph{co-nondeterministic} measure; (b) a superpolynomial lower bound for
\emph{every} $R$ is equivalent to $\NP\ne\coNP$, strictly the stronger
conjecture; (c) a lower bound for one fixed $R$ says nothing about $\Pc$ versus
$\NP$ unless paired with $\mathrm{Bridge}(R)$. So the bridge carries all
the deterministic content.

\begin{proposition}[PROVED/EXACT; adopted from Codex 1203]\label{prop:equiv}
If $\Wb_R(N)\ne\mathrm{poly}(N)$ is proved, then $\mathrm{Bridge}(R)\iff \SAT\notin\Pc$.
\end{proposition}
\begin{proof}
$\mathrm{Bridge}(R)\wedge\neg\mathrm{poly}$ gives $\SAT\notin\Pc$ by modus
tollens. Conversely $\SAT\notin\Pc$ makes the antecedent false.
\end{proof}

\section{Local multisign systems are simulated by resolution}
\begin{lemma}[PROVED/EXACT]\label{lem:local}
Let $R$ be local with parameters $(k,w)$. Every $R$-refutation of $F$ with
$L$ lines yields a resolution refutation of $F$ with at most
$L\cdot 3^{kw}$ lines. Hence $\Wb_{\mathrm{Res}}(N)\le 3^{kw}\,\Wb_R(N)$.
\end{lemma}
\begin{proof}
Process the $R$-lines in order, keeping the invariant that for every line $C$
already processed, the resolution proof contains some $D_C\subseteq C$.
Initial clauses satisfy the invariant with $D_C=C$. Let $C$ be derived soundly from
$C_1,\dots,C_j$ with $j\le k$. Then $D_{C_1},\dots,D_{C_j}\models C_1\wedge\dots\wedge C_j\models C$.
Let $X$ be the at most $kw$ variables occurring in $D_{C_1},\dots,D_{C_j}$.
Clauses over the variables of $C$ outside $X$ cannot help, so the
implicational completeness of resolution, applied to these premises restricted to $X$, gives a derivation
of some $D\subseteq C$ using only clauses over $X$. There are at most $3^{|X|}$ distinct
such clauses, so the derivation adds at most $3^{kw}$ lines. (Explicitly: either
$D_{C_1}\wedge\dots\wedge D_{C_j}$ is unsatisfiable, and resolution derives $\emptyset$,
or each implied clause on $X$ has a resolvent-closure subclause.)
For $C=\bot$ the final $D$ is $\emptyset$.
\end{proof}
\noindent COMPUTATIONAL check: the closure claim (some $D\subseteq C$ for
every implied $C$, closure size $\le3^{|X|}$) was checked on 300 random
premise tuples ($k\le3$, $w\le3$, 6 variables), covering 10{,}583 implied clauses. There were 0
failures; the largest closure had 10 clauses against a bound of 729.
Script \texttt{research/agents/claude/locality-dichotomy/local\_rule\_resolution\_sim.py},
data \texttt{\ldots/local\_rule\_resolution\_sim\_results.json}, runtime 1.4\,s, one chunk.

\begin{corollary}[PROVED/EXACT given known theorems]
For every local $R$: $\Wb_R(N)\ge 2^{\Omega(N)}$ (via Tseitin on bounded-degree expanders or random 3-CNF,
$N$ the formula size). By Proposition~\ref{prop:equiv},
$\mathrm{Bridge}(R)\iff\Pc\ne\NP$.
\end{corollary}

So for a multisign hypergraph reducer \emph{without} new nodes and with
bounded hyperedge width, the lower-bound half is already a theorem and
\emph{all} of the difficulty sits in the bridge, which is exactly $\Pc\ne\NP$.
Splitting the problem into these two halves gives no reduction in difficulty.
The same holds for any $R$ p-simulated by a system with a known
lower bound: tree-like/regular/general resolution, cutting planes,
polynomial calculus, bounded-depth Frege, AC$^0[p]$-Frege for the Lu--Santhanam--Tzameret family.

\section{The extended branch: the bridge is a feasible-soundness principle}
\begin{theorem}[Locality dichotomy; PROVED/EXACT as a case split]\label{thm:dichotomy}
For a multisign reduction system $R$, exactly one of the following holds, according to current knowledge:
\begin{enumerate}
\item[(L)] $R$ is p-simulated by a system with a proved superpolynomial lower bound. Then $\mathrm{Bridge}(R)\iff\Pc\ne\NP$.
\item[(E)] No such lower bound is known; this includes every system that p-simulates Frege or EF, and so every system with auxiliary-node extension. Then the lower-bound half is itself open. For EF it is a major open problem; no superpolynomial lower bound is known even for Frege.
\end{enumerate}
\end{theorem}
Case (E) is where the program could be nontrivial. Its bridge is the following.

\begin{proposition}[KNOWN THEOREM transfer: Cook 1975; Krajíček--Pudlák 1989]\label{prop:refl}
For a polynomial-time algorithm $A$ that decides SAT, let $S_A$ be the system in which $F$
is refuted by the empty proof whenever $A(F)=0$. Let $\mathrm{RFN}_{A,n}$
be the tautologies $[A(F)=0]\to\neg F(a)$ for $|F|=n$, written as polynomial-size
circuits with $F$ and $a$ as free bits. Then:
\begin{enumerate}
\item If $\mathrm{EF}$ has $\mathrm{poly}(n)$-size proofs of $\mathrm{RFN}_{A,n}$, then $\Wb_{\mathrm{EF}}(N)\le\mathrm{poly}(N)$.
    Proof: on input $F$, EF evaluates $A$'s circuit on the constants $F$ in polynomial size, substitutes, and applies modus ponens.
\item If $\mathrm{PV}\vdash\forall F,a\,(A(F)=0\to\neg\mathrm{Sat}(F,a))$, then (1) holds (Cook's translation).
\item Hence $\mathrm{Bridge}(\mathrm{EF})\iff\big[\SAT\in\Pc\Rightarrow\exists$ such an $A$ with polynomial-size EF proofs of $\mathrm{RFN}_{A}\big]$.
The direction $\Leftarrow$ is (1). For $\Rightarrow$, a p-bounded EF proves every tautology, $\mathrm{RFN}_A$ included, in polynomial size.
\end{enumerate}
\end{proposition}
So $\mathrm{Bridge}(\mathrm{EF})$ says: \emph{if a correct polynomial-time SAT algorithm
exists, the soundness of some such algorithm has short EF proofs.} This is a
uniform-to-feasible correctness principle. Truth of $\SAT\in\Pc$ does not
force it: correctness of $A$ need not be provable in PV or have short EF proofs.
The bridge is exactly the missing guarantee.

\paragraph{What the program yields if the bridge stays OPEN.}
A superpolynomial lower bound on $\Wb_{\mathrm{EF}}$ alone would give, through (2),
that \emph{no} polynomial-time SAT algorithm has PV-provable soundness. Equivalently,
PV does not prove that any specific polynomial-time algorithm decides SAT. This is a
strong metamathematical theorem and a known consequence; see Krajíček's book and
arXiv:1905.12935. It is not $\Pc\ne\NP$.

\section*{Barrier-transfer and PQ audit}
\textbf{Directly transferable theorem:} implicational completeness, plus the
counting of clauses over $X$ (Lemma~\ref{lem:local}), and Cook--Reckhow.
\textbf{Structural analogy needing a new theorem:} VLSI/electronics locality.
Bounded fan-in, bounded-width local rewriting is simulable at constant overhead, as in
Lemma~\ref{lem:local}. The power of extension variables is like adding
new wires or signals, which is where circuit lower bounds also stall.
\textbf{Failed transfer:} treating hard-for-resolution families
(PHP, Tseitin) as hard for the extended branch fails, since PHP has polynomial-size Frege proofs (Buss).
\textbf{PQ ledger:} item 4 (collision entropy) and item 6 (special
soundness) do not change Lemma~\ref{lem:local}. Item 5 (black-box extraction) is the
right shape for a bridge, since it extracts an EF proof from an arbitrary solver.
Proposition~\ref{prop:refl} shows that what must be extracted is a short soundness proof.
Classification: ANALOGY/CANDIDATE, and no bridge is supplied.
\textbf{Relativization:} Lemma~\ref{lem:local} and Proposition~\ref{prop:equiv} are
model-internal. Any proof of $\mathrm{Bridge}(R)$ for a case-(L) system proves $\Pc\ne\NP$,
so it inherits the relativization, natural-proofs and algebrization barriers in full.

\section*{What changed relative to root main}
Nothing in the ARC/FLMV spine changes. This run adds a reusable classification for the user's
$\Wb_R$ directive. Local reducers fall in case (L), where the bridge is the whole problem.
Extended reducers fall in case (E), where the bridge is a feasible-soundness principle and the lower bound is open.
Root \texttt{main.tex} was not edited: the Claude routine has no ownership region.

\section*{Adopted / rejected / retired prior components}
Adopted: the Codex 1203 equivalence (Proposition~\ref{prop:equiv}).
Rejected as a route: any plan to prove $\Pc\ne\NP$ by combining a resolution-type
lower bound with a ``weaker'' bridge. By Theorem~\ref{thm:dichotomy}(L), that bridge
is not weaker. Retired: none.

\section*{Stale or overstated claims found}
None in the files reviewed. Caution for future runs: the phrase ``weakest sufficient bridge'' should not be read as
``easier than $\Pc\ne\NP$''. For case (L) systems the two are provably equivalent.

\section*{Status table}
\begin{longtable}{p{0.55\textwidth}p{0.35\textwidth}}
\toprule
Claim & Status\\ \midrule
Free-sequence $\Wb_R$ poly for some $R$ $\iff\NP=\coNP$ & KNOWN THEOREM\\
Lemma \ref{lem:local} (local $\to$ resolution, factor $3^{kw}$) & PROVED/EXACT\\
Closure check, 10{,}583 implied clauses & COMPUTATIONAL\\
Bridge$(R)\iff\Pc\ne\NP$ for case (L) & PROVED/EXACT\\
Bridge(EF) $\iff$ short EF soundness proofs under $\SAT\in\Pc$ & KNOWN THEOREM transfer\\
Bridge(EF) & OPEN\\
Superpolynomial $\Wb_{\mathrm{EF}}$ / $\Wb_{\mathrm{Frege}}$ & OPEN\\
P versus NP & OPEN\\ \bottomrule
\end{longtable}

\section*{Single next target}
Within case (E), the most useful falsifiable sub-target is the following. Fix an explicit polynomial-time
algorithm $A$ that is \emph{correct on a restricted class} $\mathcal C$, for example
2-SAT, Horn-SAT, XOR-SAT, or bounded-treewidth CNF. Determine the EF/Frege size of $\mathrm{RFN}_{A}$
restricted to $\mathcal C$. Short proofs are expected and would confirm the mechanism.
A provably \emph{superpolynomial} gap on one natural class, between deciding by $A$ and proving
$A$'s soundness, would be the first evidence that the feasible-soundness principle can fail.
It would also show that the bridge cannot be proved from correctness alone.
For XOR-SAT the relevant reflection carries Gaussian-elimination soundness. That is known to be hard for
bounded-depth Frege (Tseitin) but polynomial in Frege. This makes XOR-SAT a calibrated test case.

% END_OF_VERIFIED_REPORT::CLAUDE::WORSTCASE-BRIDGE-0001
\end{document}
